% Transcription — fonds Alexandre Grothendieck, shelfmark 156-9, batch 4 (pages 61-80).
% Established from the facsimile archives/batches/156-9/batch-04.pdf (a mirror of
% https://grothendieck.umontpellier.fr/156-9.pdf), per the skill
% transcribe-grothendieck.
%
% Pass: Opus 5.5 (claude-opus-5-5), 2026-10-03 — first pass, unchecked against the pages by a human.
%
% All twenty pages are transcribed. Each sheet carries his own pagination
% (58-77): the archivists' 61 is his 58.
% It is recorded once per page in a \note, because the pages refer to each
% other by his numbers (« p. 59 »).
%
% Content: « dispositions » (a set with a symmetric relation |o|), the
% operators supp° and cosupp° on its power set, the lattice Σ_E, the attempt
% to define correspondences between dispositions and their transposes; then,
% from the 10 July restart (p. 75), a fresh write-up: dispositions, their
% relation to strict graphs, and the notion of « spatiale » (complete lattice
% with an order-reversing involution, Spat 1-3), Proposition 1 and the
% beginning of Proposition 2, which runs on past the batch.
%
% Notation fixed for the batch: the operators he writes « s--pp° »,
% « cos--pp° » are set \operatorname{supp}^{\circ}, \operatorname{cosupp}^{\circ};
% his big « C » before a set or an element of a spatiale is \complement;
% the disposition relation is \mathrel{|o|}; \bar\alpha is the extension of a
% correspondence α to parts.
%
% The dating is the inventory's own for the shelfmark; the title comes from
% src/content/catalogue.ts verbatim, its bracketed words supplied by the
% archivists.
\documentclass[11pt,a4paper]{article}
\input{../preamble/grothendieck}

\folder{156-9}
\batch{4}
\pages{61}{80}
\foldertitle{[Chapitre] IX et IX bis. [Ateliers] : notes manuscrites (05-15/07/1986).}
\dating{1986}
\watermark{Édition de démonstration}

\begin{document}

\note{l'argument vient d'avant ce lot~: la page~61 poursuit une discussion, commencée plus haut, de la relation de disposition $|o|$ et des opérateurs $\operatorname{supp}^{\circ}$, $\operatorname{cosupp}^{\circ}$, dont les définitions ne figurent pas ici.}

\page{61}\note{p. 58 de l'auteur}
On donne une catégorie, \struck{la catégorie} \uncertain{telle}, et un objet (équivalent~: un graphe (strict)), \uncertain{constituant} un \ill{}, exprimant l'intuition qui s'attache à la situation. On pourrait appeler un $(E, |o|_E)$ un ensemble muni d'une relation de disposition, ou \emph{disposition}, et \struck{la catégorie précédente} parler de \emph{morphismes de dispositions}, notation $\operatorname{Hom}_{\mathrm{dis}}(\underline{E}, \underline{E}')$, et la \emph{catégorie des dispositions}. Ceci dit, \struck{\ill{}} \add{tel} pour un morphisme de dispositions
\[ E \xrightarrow{\ \alpha\ } \mathfrak{P}(E') \]
considérons l'application
\[ \bar\alpha : \mathfrak{P}(E) \longrightarrow \mathfrak{P}(E'), \qquad A \longmapsto \bigcup_{x \in A} \alpha(x), \]
et l'application composée
\[ \Sigma_E \xhookrightarrow{\ \mathrm{inc}\ } \mathfrak{P}(E) \xrightarrow{\ \bar\alpha\ } \mathfrak{P}(E') \xrightarrow{\ \operatorname{supp}^{\circ}_{E'}\ } \Sigma_{E'} ; \]
\note{l'opérateur noté ici $\operatorname{supp}^{\circ}$ (et plus loin $\operatorname{cosupp}^{\circ}$) est abrégé « s--pp° », « cos--pp° » sur les pages~; la lecture « supp » est probable mais non assurée. Sur cette page, une première écriture biffée de $\Sigma_E$ précède la composée.}
\struck{Il est.} On pose
\[ \alpha_{*} : \Sigma_E \longrightarrow \Sigma_{E'}, \qquad \alpha_{*}(S) = \operatorname{supp}^{\circ}_{E'}\bigl(\bar\alpha(S)\bigr). \]
Cette application \ill{} on fait défini\ill{} \struck{\ill{}} \add{\uncertain{cueillies}}\note{fin de ligne en bas de page, très incertaine.}

\page{62}\note{p. 59 de l'auteur}
quel que soit la correspondance $\alpha$, même si celle-ci n'est pas \emph{compatible} aux dispositions. Cependant, quand il y a compatibilité, on trouve \uncertain{aussi} la formule essentielle
\[ (*) \qquad \boxed{\ \alpha_{*}\bigl(\operatorname{supp}^{\circ}_{E} A\bigr) = \operatorname{supp}^{\circ}_{E'}\bigl(\bar\alpha(A)\bigr)\ } \]

\note{toute la démonstration qui suit, jusqu'à « Il faudrait », est barrée de longs traits obliques.}
\struck{\emph{Dém.} Soit $\bar A = \operatorname{supp}^{\circ}_{E}(A)$, donc $A \subset \bar A$~; alors}
\[ \text{\struck{$\alpha_{*}(\bar A) = \operatorname{supp}^{\circ}_{E'}(\bar\alpha(\bar A)) \supset \operatorname{supp}^{\circ}_{E'}(\bar\alpha(A))$}} \]
\struck{puisque $\bar A \supset A$ implique que $\bar\alpha(\bar A) \supset \bar\alpha(A)$. Prouvons l'inclusion inverse}
\[ \text{\struck{$\operatorname{supp}^{\circ}_{E'}(\bar\alpha(\bar A)) \subset \operatorname{supp}^{\circ}_{E'}(\bar\alpha(A))$,}} \]
\struck{donc que pour tout $x' \in E'$, on a}
\[ \text{\struck{$x' \mathrel{|o|} \bar\alpha(A) \Longrightarrow x' \mathrel{|o|} \bar\alpha(\bar A)$~?}} \]
\struck{\ill{} \ill{} \ill{}} \ill{} \ill{} \uncertain{ceci} faux~! Il faudrait \ill{}

\uncertain{Mais il existe une} \ill{} \uncertain{à l'aide de la formule} \uncertain{« duale »}
\[ (**) \qquad \alpha_{*}\bigl(\operatorname{cosupp}^{\circ}_{E}(B)\bigr) = \operatorname{cosupp}^{\circ}_{E'}\bigl(\bar\alpha(B)\bigr) \]
\struck{\ill{}} $(\forall B \in \mathfrak{P}(E))$ \ill{} \uncertain{qu'on applique} (\uncertain{une fois pour $B$, une fois pour $A$}) \add{deux fois} \struck{\ill{} fois} \uncertain{dans le cas où} $B = \operatorname{cosupp}^{\circ}_{E}(A)$ \struck{\ill{} fois pour $A$}~:
\begin{align*}
\alpha_{*}\bigl(\operatorname{supp}^{\circ}(B)\bigr) &= \operatorname{cosupp}^{\circ}_{E'}\bigl(\bar\alpha(\operatorname{cosupp}^{\circ}_{E}(A))\bigr) \\
&= \operatorname{cosupp}^{\circ}_{E'}\Bigl(\operatorname{cosupp}^{\circ}_{E'}\bigl(\bar\alpha(\ \text{---}\ )\bigr)\Bigr)
\end{align*}
\note{sous la seconde ligne, une accolade porte $\alpha_{*}(\operatorname{cosupp}^{\circ}(A))$, suivi d'un « $= \operatorname{cosupp}^{\circ}$ » biffé~; les indices et la notation exacte de cette fin de page sont incertains.}

\page{63}\note{p. 60 de l'auteur}
et comme $\alpha_{*}\operatorname{cosupp}^{\circ}_{E}(A) = \operatorname{cosupp}^{\circ}_{E'}(\bar\alpha(A))$, on trouve finalement
\begin{align*}
\uncertain{dans}\ \operatorname{supp}^{\circ}_{E}(A) &= \operatorname{cosupp}^{\circ}_{E'}\bigl(\operatorname{cosupp}^{\circ}_{E'}(\bar\alpha(A))\bigr) \\
&= \operatorname{supp}^{\circ}_{E'}(\bar\alpha(A)) \qquad \text{cqfd.}
\end{align*}
\note{le membre de gauche de la première ligne est tel qu'écrit~; on attendrait $\alpha_{*}(\operatorname{supp}^{\circ}_{E}(A))$.}

Donc il faut établir la formule, qui, pour $B \in \Sigma_E$, \struck{donne} s'écrit
\begin{align*}
\alpha_{*}(\complement B) &= \operatorname{cosupp}^{\circ}_{E'}\bigl(\bar\alpha(B)\bigr) \\
&= \operatorname{cosupp}^{\circ}_{E'}\bigl(\underbrace{\operatorname{supp}^{\circ}_{E'}(\bar\alpha(B))}_{\overset{\text{déf}}{=}\ \alpha_{*}(B)}\bigr) = \complement\,\alpha_{*}(B),
\end{align*}
i.e.
\[ \boxed{\ \alpha_{*}(\complement B) = \complement\,\alpha_{*}(B)\ } \]
i.e.\ $\alpha_{*}$ compatible avec $\complement$.\note{le symbole lu $\complement$ est un grand « C » placé devant la partie, qu'on lit comme le complémentaire.}

\struck{Soit donc} Il faut donc prouver $(**)$, qui apparaît comme la formule-clef. Posons
\[ C = \operatorname{cosupp}^{\circ} B = \{\, x \in E \mid x \mathrel{|o|} B \,\}~; \]
alors la formule $(**)$ s'écrit
\[ \operatorname{supp}^{\circ}_{E'}\bigl(\bar\alpha(C)\bigr) = \operatorname{cosupp}^{\circ}_{E'}\bigl(\bar\alpha(B)\bigr) \qquad \bigl(C = \operatorname{cosupp}^{\circ}_{E}(B)\bigr). \]
Prouvons d'abord $\subset$, qui revient à
\[ \bar\alpha(C) \subset \operatorname{cosupp}^{\circ}_{E'}\bigl(\bar\alpha(B)\bigr), \]
ou encore
\[ \bar\alpha(C) \mathrel{|o|_{E'}} \bar\alpha(B) \]
\struck{$B$ \ill{}} et résulte de

\page{64}\note{p. 61 de l'auteur}
\[ B \mathrel{|o|_{E}} C, \]
\uncertain{compte tenu} de la \struck{propriété fondamentale} \add{compatibilité} de $\alpha$ aux dispositions. Reste la question
\[ \operatorname{cosupp}^{\circ}_{E'}\bigl(\bar\alpha(B)\bigr) \subset \operatorname{supp}^{\circ}_{E'}\bigl(\bar\alpha(C)\bigr), \]
i.e.\ ceci~:

$(***)$ Soit $x' \in E'$, alors $x' \mathrel{|o|} \bar\alpha(B)$ implique
\[ x' \in \operatorname{supp}^{\circ}_{E'}\bigl(\bar\alpha(C)\bigr), \]
i.e.\ $\forall y' \in E'$, avec $y' \mathrel{|o|} \bar\alpha(C)$, on a $x' \mathrel{|o|} y'$.

Pour le dire autrement~:
\[ x', y' \in E', \quad x' \mathrel{|o|} \bar\alpha(B), \quad y' \mathrel{|o|} \bar\alpha(C) \ \Longrightarrow\ x' \mathrel{|o|} y'. \]

\marginal{Mais, en effet, on a plus \ill{}}
Je ne vois pas que cela puisse se déduire de la seule compatibilité de $\alpha$ aux dispositions, qui affirme \add{seulement} que deux \uncertain{éléments} de $E'$ \struck{\ill{}} sont \uncertain{séparés} \add{\uncertain{dans} la disposition} \uncertain{d'éléments} \add{$x', y'$} \ill{} de $E'$ qui \struck{\ill{}} \uncertain{correspondent} \ill{}~: des éléments $x, y$ de $E$. \struck{Il semble donc, pour pouvoir}

\struck{\uncertain{conclure}} D'ailleurs, \ill{} un contre-exemple immédiat, en prenant \uncertain{dans} $E$, $E'$ \add{\ill{}} la « disposition \uncertain{tautologique} » $x \mathrel{|o|} y \Longleftrightarrow x \neq y$, d'où $\Sigma_E = \mathfrak{P}(E)$, $\Sigma_{E'} = \mathfrak{P}(E')$, et une \emph{application} $\alpha : E \to E'$, d'où $\alpha_{*} : \Sigma_E \to \Sigma_{E'}$ \uncertain{donnée par} $A \mapsto \alpha(A)$. \struck{\emph{Mais} \ill{}, \uncertain{pour} cette application} \add{la} compatibilité de $\alpha$ aux dispositions signifie \uncertain{ici} $\alpha$ injection, mais $\alpha_{*}$ n'est \uncertain{guère} compatible~: $\complement$

\page{65}\note{p. 62 de l'auteur}
que si de plus $\alpha$ est surjection, i.e.\ bijection~!

Donc il faut une condition supplémentaire, pour pouvoir affirmer les commutations de $\alpha_{*}$ à $\complement$, ou mieux, la formule $(**)$. On suppose donnée une correspondance en sens inverse
\[ E' \xrightarrow{\ \beta\ } \mathfrak{P}(E), \]
i.e.\ une application
\[ \mathfrak{P}(E') \xrightarrow{\ \bar\beta\ } \mathfrak{P}(E) \]
commutant aux réunions. On \uncertain{suppose} que
\[ (*) \qquad \forall x' \in E', \quad \bar\alpha\bigl(\beta(x')\bigr) \subset \operatorname{supp}^{\circ}_{E'}(\{x'\}). \]
Si $B \subset E$, \struck{alors} et $x' \in E'$, alors
\[ x' \mathrel{|o|} \bar\alpha(B) \Longleftrightarrow \bar\alpha\bigl(\beta(x')\bigr) \mathrel{|o|} \bar\alpha(B) \]
et, si $\bar\alpha$ a une propriété de \emph{fidélité} pour $|o|$, on en conclut
\[ \beta(x') \mathrel{|o|} B. \]
\marginal{\textbf{NB} $(*)$ est satisfait si \ill{} \uncertain{l'itération} \ill{}, \uncertain{cf.\ p.~24} \ill{}}
Appliquons ceci~; à la situation $(**)$ (p.\ précédente)~: $(B, x')$ et~: $(C, y')$, on trouve
\[ \beta(x') \mathrel{|o|} B, \qquad \underbrace{\beta(y') \mathrel{|o|} C = \operatorname{cosupp}^{\circ}(B)}_{\text{i.e. } \beta(y') \subset \operatorname{supp}^{\circ}(B)} \]
d'où
\[ \beta(x') \mathrel{|o|} \beta(y'), \qquad \text{d'où} \]
\note{la référence « p.\ précédente » renvoie, selon la numérotation de l'auteur, à la formule $(**)$ de la p.~59 (ici p.~62 des archivistes).}

\page{66}\note{p. 63 de l'auteur}
\[ \bar\alpha\bigl(\beta(x')\bigr) \mathrel{|o|} \bar\alpha\bigl(\beta(y')\bigr) \]
(compatibilité de $\alpha$ aux dispositions), pour en conclure, comme on voudrait, $x' \mathrel{|o|} y'$, il faudrait que l'inclusion $(*)$ soit une identité, ou plutôt que
\[ \operatorname{supp}^{\circ}_{E'}\Bigl(\bar\alpha\bigl(\beta(x')\bigr)\Bigr) = \operatorname{supp}^{\circ}_{E'}\bigl(\{x'\}\bigr). \]

Voici un choix \emph{optimal} de $\beta$, en termes de $\alpha$~:
\begin{gather*}
\text{Soit } \beta = {}^{t}\alpha^{*} \text{ défini par} \\
\alpha^{*}(x') = \bigl\{\, x \in E \mid \alpha(x) \subset \operatorname{supp}^{\circ}_{E'}(\{x'\}) \,\bigr\}
\end{gather*}
\struck{Notre hypothèse supplémentaire sur \uncertain{ce choix} \ill{} $\beta$ \ill{}, \ill{} On \ill{}} \struck{Avec cette \ill{}} \struck{Avec cette}
donc tautologiquement
\[ \bar\alpha\bigl(\alpha^{*}(x')\bigr) \subset \operatorname{supp}^{\circ}_{E'}(\{x'\}), \]
\struck{et il faut supposer que cette inclusion} \struck{\uncertain{définit}} \add{\uncertain{donne}} et il faut supposer que cette inclusion est une \emph{égalité}~:
\[ \text{(i)} \qquad \boxed{\ \operatorname{supp}^{\circ}_{E'}\Bigl(\bar\alpha\bigl(\alpha^{*}(x')\bigr)\Bigr) = \operatorname{supp}^{\circ}_{E'}(\{x'\})\ } \]
(remplace l'hypothèse de surjectivité dans l'exemple \uncertain{trivial}). De plus, la propriété que $\alpha$ soit compatible, doit être remplacée par une hypothèse de « compatibilité fidèle »
\[ \text{(ii)} \qquad \boxed{\ \text{si } x, y \in E,\quad x \mathrel{|o|} y \Longleftrightarrow \bar\alpha(x) \mathrel{|o|} \bar\alpha(y)\ } \]
\note{dans (i), le membre de gauche porte « $\operatorname{supp}^{\circ}$ » et l'indice $E'$ est à peine visible~; dans (ii), $\bar\alpha(x)$ s'entend $\alpha(x)$.}

\page{67}\note{p. 64 de l'auteur}
qui implique aussitôt, pour $A, B \subset E$
\[ \bar\alpha(A) \mathrel{|o|} \bar\alpha(B) \Longleftrightarrow A \mathrel{|o|} B. \]

Ces conditions deviennent \uncertain{moins} draconiennes \uncertain{donnant ce} qu'il nous faut, mais peut-être même un peu \struck{plus} trop --- impliquent-elles pas presque
\[ \alpha_{*} : \Sigma_E \xrightarrow{\ \sim\ } \Sigma_{E'} \]
bijectif~? \struck{\ill{}} D'autre part, on aimerait peut-être \struck{mais} avoir la formule $(*)$ p.~59, \uncertain{dans} des cas où \ill{} la formule $(**)$ p.~59 est fausse (comme on \ill{} dans l'exemple des \uncertain{tautologies}). Donc je n'ai pas l'impression encore d'avoir vraiment bien compris la situation.
\note{« p.~59 » est la pagination de l'auteur~: c'est la p.~62 (numérotation des archivistes) ici.}

\emph{Dans} la formule $(*)$ p.~59
\[ \underbrace{\alpha_{*}\operatorname{supp}^{\circ}_{E}(A)}_{\overset{\text{déf}}{=}\ \operatorname{supp}^{\circ}_{E'}\bar\alpha(\operatorname{supp}^{\circ}_{E}(A))} = \operatorname{supp}^{\circ}_{E'}\bigl(\bar\alpha(A)\bigr) \]
l'inclusion $\supset$ est claire, puisque $\operatorname{supp}^{\circ}_{E}(A) \supset$ \struck{$\operatorname{supp}^{\circ}_{E}$} $A$ d'où
\[ \bar\alpha\bigl(\operatorname{supp}^{\circ}_{E}(A)\bigr) \supset \bar\alpha(A) \]
et il suffit de passer aux $\operatorname{supp}^{\circ}_{E'}$. Donc le pb est

\page{68}\note{p. 65 de l'auteur}
\[ \underbrace{\alpha_{*}\bigl(\operatorname{supp}^{\circ}_{E}(A)\bigr)}_{\operatorname{supp}^{\circ}_{E'}(\bar\alpha(\operatorname{supp}^{\circ}_{E}(A)))} \overset{?}{\subset} \operatorname{supp}^{\circ}_{E'}\bigl(\bar\alpha(A)\bigr) \]
ou encore dans ceci~:
\begin{align*}
&\text{Si } B \subset \operatorname{supp}^{\circ}_{E}(A) \text{ i.e. } \forall x \in E,\ x \mathrel{|o|} A \Rightarrow x \mathrel{|o|} B, \\
&\text{peut-on en conclure que} \\
&\overset{?}{\Bigl\{}\ \bar\alpha(B) \subset \text{\struck{$\bar\alpha(\operatorname{supp}$}}\ \operatorname{supp}^{\circ}_{E'}\bigl(\bar\alpha(A)\bigr) \text{ i.e. } \forall x' \in E', \\
&\qquad x' \mathrel{|o|} \bar\alpha(A) \Rightarrow x' \mathrel{|o|} \bar\alpha(B)~?
\end{align*}

Pour l'affirmer, il faudrait une \uncertain{cor\-res\-pon\-dance} en sens inverse
\[ \beta : E' \longrightarrow \mathfrak{P}(E) \]
telle que l'on ait, $\forall x \in E$, $x' \in E'$
\begin{align*}
(*) \qquad & x \mathrel{|o|} \beta(x') \Longleftrightarrow \alpha(x) \mathrel{|o|} x', \\
&\text{d'où, pour } A \subset E,\ B' \subset E' \\
(**) \qquad & A \mathrel{|o|} \bar\beta(B') \Longleftrightarrow \bar\alpha(A) \mathrel{|o|} B'
\end{align*}
\struck{Dans ce cas}, l'implication \uncertain{ci-dessus} \uncertain{se} \struck{\ill{}} \uncertain{reformule} \ill{}
\[ \text{\struck{$x \mathrel{|o|} \beta(x) \Rightarrow$}} \qquad \beta(x') \mathrel{|o|} A \overset{?}{\Longrightarrow} \beta(x') \mathrel{|o|} B \]
ce qui résulte bien de $B \subset \operatorname{supp}^{\circ}_{E}(A)$.

\marginal{On \uncertain{n'a} \uncertain{plus} \ill{} $\forall x', \ldots$ \ill{} $x \mathrel{|o|} y \Rightarrow \ldots$~!}
Ici encore, s'il existe de telles corr.\ $\beta$ \uncertain{cachées}, il doit y en avoir une optimale~: soit en effet
\[ \gamma(x') = \{\, x \in E \mid \alpha(x) \mathrel{|o|} x' \,\}, \]
\uncertain{cette $\gamma$ est définie en termes de $\alpha$ seule}, alors la propriété $(*)$ sur $\beta$ se lit
\[ x \mathrel{|o|} \beta(x') \Longleftrightarrow x \in \gamma(x') \qquad \text{i.e.} \]
\[ \gamma(x') = \operatorname{cosupp}^{\circ}\bigl(\beta(x')\bigr) \qquad \text{\struck{\ill{}}} \]
\note{la note marginale, écrite en biais à gauche du bas de page, est presque entièrement illisible.}

\page{69}\note{p. 66 de l'auteur}
ce qui signifie deux choses~:
\begin{align*}
1^{\circ})\quad & \gamma(x') \in \Sigma_E \\
2^{\circ})\quad & \text{\struck{$\beta$}}\ \operatorname{supp}^{\circ}\beta(x') = \complement\,\gamma(x') \quad \bigl(= \operatorname{cosupp}^{\circ}(\gamma(x'))\bigr)
\end{align*}
Donc une telle $\beta$ existe ssi \ill{}, et on peut prendre
\[ \beta(x') = \operatorname{cosupp}^{\circ}\bigl(\gamma(x')\bigr) = \bigl\{\, x \in E \mid \forall y \in E,\ \alpha(y) \mathrel{|o|} x' \Longrightarrow x \mathrel{|o|} y \,\bigr\}~; \]
cette $\beta$ est sûrement la « plus grande » des \uncertain{telles} \add{\uncertain{correspondances}}, et \ill{} \uncertain{satisfaisant} $(*)$ [et en \uncertain{particulier}~: \uncertain{étant} donnée une $\beta$ quelconque, on a \ill{} \uncertain{compte tenu} de \uncertain{la} \uncertain{fonction} $\operatorname{supp}^{\circ}_{E}$]. On notera d'ailleurs que \struck{cette} $\beta$ les $\beta$ associées à $\alpha$ ne changent pas, quand on remplace $\alpha$ par son composé avec $\operatorname{supp}^{\circ}_{E'}$.

\marginal{\textbf{NB} Si $E, E'$ sont les dispositions \uncertain{tautologiques} ($x \mathrel{|o|} y \Leftrightarrow x \neq y$), alors on voit \ill{} $\beta$ \ill{} $\operatorname{supp}^{\circ} = \mathrm{id}$ \ill{} \ill{} \ill{} \ill{}}
\note{note marginale de neuf lignes en biais, à gauche du milieu de page~; seuls quelques mots et formules s'en lisent.}

Je vois ici \uncertain{déjà} \uncertain{poindre} une proposition.

\note{ce qui suit est barré de deux longs traits obliques.}
\struck{\emph{Proposition}. Soient $(E, |o|_E)$, $(E', |o|_{E'})$ deux dispositions.} \struck{Il y a} \struck{Soit $\operatorname{Corr}(E, E')$ l'ensemble \ill{} des correspondances de $E$ dans $E'$, et $\operatorname{Corrdis}(E, E')$ le sous-ensemble \uncertain{formé} des}

\page{70}\note{p. 67 de l'auteur}
\note{le haut de la page, jusqu'à « ou s'il existe une », est barré de trois traits obliques légers, en plus des biffures horizontales notées.}
\struck{ce qui signifie et}

\struck{\emph{Définition}. Un couple $(\alpha_{*}, \beta_{*})$ d'applications}
\[ \text{\struck{$\alpha_{*} : \Sigma_E \longrightarrow \Sigma_{E'}$}} \]
\struck{est dit une \emph{correspondance de disposition} de $E$ à $E'$, s'il} \struck{existe \ill{} \ill{} applications} \struck{$\beta_{*} : \Sigma_{E'}$} \struck{\ill{} \uncertain{Sup} quelconques, et s'il existe une \ill{} application \ill{}}

\emph{Soit} une correspondance $\alpha$ entre deux ensembles $E, E'$ (i.e.\ une \emph{relation} entre $E, E'$ i.e.\ une partie $\alpha \subset E \times E'$) \uncertain{peut être identifiée} \uncertain{à} une application
\[ \alpha_{*} : \mathfrak{P}(E) \longrightarrow \mathfrak{P}(E') \]
\uncertain{commutant} aux $\bigcup$ quelconques.
\marginal{Correspondance \ill{} $\Longleftrightarrow$ \ill{} \ill{} $\mathfrak{P}(E)$.}

De même, \struck{\ill{}} si $E, E'$ sont des dispositions, \uncertain{on} considère les applications
\[ \alpha_{*} : \Sigma_E \longrightarrow \Sigma_{E'} \]
commutant aux Sup quelconques, qui généralisent \uncertain{ainsi} les correspondances des dispositions « triviales ». Je suppose \uncertain{donnée} l'application
\[ \alpha : E \longrightarrow \Sigma_{E'}, \qquad x \longmapsto \alpha_{*}\bigl(\underbrace{\operatorname{supp}^{\circ}_{E}(\{x\})}_{s(x)}\bigr) \]

\page{71}\note{p. 68 de l'auteur}
et on \struck{trouve ainsi} l'étend en $\alpha_{*}$ par
\[ \alpha_{*}(S) = \mathop{\mathrm{Sup}}\limits_{x \in S\ (\text{dans } \Sigma_{E'})} \bigl(\underbrace{\alpha(x)}_{\alpha_{*}(s(x))}\bigr) = \operatorname{supp}^{\circ}_{E'}\bigl(\bar\alpha(S)\bigr) \]
\emph{forcément}, puisque
\[ S = \Bigl[\mathop{\mathrm{Sup}}\limits_{x \in S\ (\text{dans } \Sigma_E)}\Bigr](s_x). \]
Ainsi

\emph{Lemme}. Corr.\ bijective entre \struck{Appl} $\operatorname{Hom}_{\mathrm{ens}}(E, \Sigma_{E'})$ et $\operatorname{Hom}_{\mathrm{ord}}^{!}(\Sigma_E, \Sigma_{E'})$, où le ! signifie \uncertain{qu'on} prend les hom.\ d'ens.\ ordonnés \uncertain{qui} \struck{\ill{}} \uncertain{commutent} aux Sup quelconques.

\marginal{Mais il n'est pas clair \uncertain{du tout} que si $A \subset E$, \uncertain{on ait} $\alpha_{*}(\operatorname{supp}^{\circ}_{E}(A)) \overset{?}{=} \operatorname{supp}^{\circ}_{E'}(\alpha_{*}(A))$, \ill{} \ill{}}
On voudrait cependant une notion de \uncertain{transposition} d'un tel \uncertain{morphisme}, qui généralise le cas des « dispositions triviales », où la définition des transposés $\beta$ de $\alpha$ peut se décrire par
\[ \text{\struck{$\alpha(x) \ni$}}\qquad x' \in \alpha(x) \Longleftrightarrow x \in \beta(x') \]
\uncertain{ici} on pose~: la \emph{négation}
\[ x' \mathrel{|o|} \alpha(x) \Longleftrightarrow x \mathrel{|o|} \beta(x') \]
\uncertain{Ou alors}
\note{la note marginale, écrite en biais sur la gauche, est en partie illisible~; le dernier mot de la page reste en suspens.}

\page{72}\note{p. 69 de l'auteur}
\emph{Lemme 2}. Considérons une application
\[ \alpha_{*} : \Sigma^{*}_{A} \longrightarrow \Sigma^{*}_{A'} \]
\struck{« opposée » à \ill{} \ill{}} \add{\uncertain{correspondant} \ill{}} \struck{\uncertain{correspondant} \ill{}}
\note{dans cette page les indices $E$, $E'$ de $\Sigma$ ont été surchargés en $A$, $A'$ (et coiffés d'une astérisque) dans les premières formules~; plus bas l'auteur revient à $E$, $E'$. On garde ce que porte chaque formule.}
donnée par le lemme 1, \uncertain{à partir} d'une
\[ \alpha : E \longrightarrow \mathfrak{P}(E'), \]
ou son composé avec $\operatorname{supp}^{\circ}_{E'}$, donc
\[ \text{\struck{i.e.\ \ill{}}}\qquad \alpha_{*}(S) = \operatorname{supp}^{\circ}_{E'}\bigl(\bar\alpha(S)\bigr). \]
Conditions \struck{équivalentes}~:

(i) $\exists\, \beta_{*} : \Sigma^{*}_{A'} \longrightarrow \Sigma^{*}_{A}$, commutant aux Sup quelconques (i.e.\ comme $\alpha_{*}$) satisfaisant, pour \struck{\ill{}} $S \in \Sigma^{*}_{A}$, $S' \in \Sigma^{*}_{A'}$
\[ (*) \qquad \alpha_{*}(S) \mathrel{|o|} S' \Longleftrightarrow S \mathrel{|o|} \beta_{*}(S') \]

(ii) $\exists\, \beta : E' \longrightarrow$ \struck{$\Sigma^{*}_{E}$} $\mathfrak{P}(E)$, telle que l'on ait, pour $x \in E$, $x' \in E'$
\[ (**) \qquad \alpha(x) \mathrel{|o|} x' \Longleftrightarrow x \mathrel{|o|} \beta(x'). \]

De plus, les applications $\beta_{*}$ dans (i) sont déterminées uniquement, et si $\beta$ est comme dans (ii), $\beta_{*}$ est \emph{associée} à $\beta$
\[ \beta_{*}(S') = \operatorname{supp}^{\circ}_{E}\bigl(\bar\beta(S')\bigr) \]
\marginal{\uncertain{Abréviations}~: \ill{} \ill{}}
\marginal{« \uncertain{Hypothèse} \uncertain{naturelle} \ill{} \ill{}~: $S \in \Sigma_{E'}$, $\{\, x \in X \mid \bar\alpha(x) \subset S \,\} \in \Sigma_E$ \ill{} \uncertain{Valide} \ill{} \uncertain{la} \uncertain{formule} $\alpha_{*}(\operatorname{supp}^{\circ}_{E}(A)) = \operatorname{supp}^{\circ}_{E'}(\bar\alpha(A))$ ».}
Enfin, les conditions (i) (ii) sont équivalentes à la suivante~:

(iii) $\forall x' \in E'$, soit
\[ \bar\beta(x') = \{\, x \in E \mid \alpha(x) \mathrel{|o|} x' \,\} \subset E~; \]
\struck{Alors la partie} on veut
\[ \bar\beta(x') \in \Sigma_E. \]
\note{les deux notes marginales, l'une au milieu, l'autre au bas de la marge gauche, sont écrites en biais et en partie illisibles.}

\page{73}\note{p. 70 de l'auteur}
Ceci posé, \struck{les} $\beta_{*}$ \uncertain{donnée} \uncertain{par}
\[ \beta_{*}(x') = \complement\bigl(\underbrace{\bar\beta(x')}_{\in\,\Sigma_E}\bigr) = \operatorname{cosupp}^{\circ}_{E}\bigl(\bar\beta(x')\bigr) \]
et l'application $\hat\beta$ dans (ii) satisfaisant la condition $(**)$ des $\beta_{*}$ est \emph{associée} à $\beta$,
\note{le symbole lu $\hat\beta$ est un $\beta$ coiffé d'un signe, de lecture incertaine.}
ou encore
\begin{align*}
\operatorname{supp}^{\circ}_{E}\beta(x') &= \complement\bigl(\bar\beta(x')\bigr) \\
\text{ou aussi}\qquad \operatorname{cosupp}^{\circ}\bigl(\beta(x')\bigr) &= \bar\beta(x').
\end{align*}

\emph{Définition}. Les $\alpha_{*}$ satisfaisant les conditions équivalentes (i) (ii) sont appelées correspondances de dispositions \struck{entre} $\underline{E}$, $\underline{E}'$. La correspondance $\beta_{*}$ (qui est aussi une corr.\ disp.), de (i), s'appelle la corr.\ disposition \struck{associée} \emph{transposée} de $\alpha_{*}$.

\emph{Scholie}. On trouve \struck{ainsi} une bijection
\[ \alpha_{*} \longmapsto {}^{t}\alpha_{*} : \operatorname{Corrdis}(\underline{E}, \underline{E}') \longrightarrow \operatorname{Corrdis}(\underline{E}', \underline{E}) \]
et \ill{}.
\[ \begin{cases} {}^{t}({}^{t}\alpha_{*}) = \alpha_{*} \\ {}^{t}({}^{t}\beta_{*}) = \beta_{*} \end{cases} \]

\page{74}\note{p. 71 de l'auteur}
De plus, l'on a \struck{alors}, pour $A \subset E$
\[ (***) \qquad \underbrace{\alpha_{*}\bigl(\operatorname{supp}^{\circ}_{E}(A)\bigr)}_{\overset{\text{déf}}{=}\ \operatorname{supp}^{\circ}_{E'}(\bar\alpha(\operatorname{supp}^{\circ}_{E}(A)))} = \operatorname{supp}^{\circ}_{E'}\bigl(\bar\alpha(A)\bigr) \]

\textbf{NB} J'ignore si la validité de cette formule, pour une \struck{$\alpha_{*}$} quasi-corr.\ dispositions donnée $\alpha_{*}$, indique que c'est une corr.\ disposition.

En analogie avec le cas des dispositions triviales, où on distingue le cas d'une correspondance qui est une application, on transposé d'une application, je dirai que $\alpha_{*}$ est une \emph{application} si ${}^{t}\alpha_{*} = \beta_{*}$ a les propriétés \emph{supplémentaires}~: il commute à $\complement$ (donc aussi aux Inf quelconques)

\page{75}\note{p. 72 de l'auteur}
\emph{10 juillet}

Je reprends dès le début la formulaire des dispositions, \uncertain{synthèse}.

1. \emph{Dispositions}, (dites \emph{relations de disposition}) \emph{simples}, \struck{\ill{}}: ensemble $E$ muni d'une relation $\mathrel{|o|}$, \struck{\uncertain{supposée} \ill{}} \uncertain{symétrique}~:

\emph{Disp.\ 1} \quad $x \mathrel{|o|} y \Rightarrow y \mathrel{|o|} x$

\emph{Disp.\ 2} \quad $x \mathrel{|o|} x \Longrightarrow x \mathrel{|o|} y\ \forall y \in E$ [« \struck{\ill{}} \uncertain{dans} \uncertain{ce} \uncertain{cas} \uncertain{isolé} » \uncertain{dans} \uncertain{la} \uncertain{disp.}] l'ens.\ $E_0$.
\marginal{\uncertain{Structures} de disposition}
On désigne par \emph{dispositions antiréflexives} \struck{\ill{}} quand $\mathrel{|o|}$ est antiréflexive (\uncertain{Disp.\ 2}), i.e.\ $E_0 = \emptyset$. Ce qui \uncertain{implique} \uncertain{Disp.\ 2}.

\emph{Relation à notion de graphe (strict)}. On associe à une str.\ de disp.\ $\mathrel{|o|}$ sur $E$, une structure de graphe, une partie \struck{\ill{}} \add{\uncertain{autre}} $A \subset \mathfrak{P}_2(E)$, l'ens.\ suivant
\[ A = \bigl\{\, \{x, y\} \in \mathfrak{P}_2(E) \mid x \mathrel{\overline{|o|}} y \,\bigr\}. \]
\note{la barre sur $|o|$ est lue comme la négation de la relation.}
On récupère (\add{presque}) la structure de dispositions à partir de cette structure de graphe, \struck{et} \add{il faut \uncertain{connaître}} \struck{\ill{}} \uncertain{plus} la partie $E_0$ de l'ens.
\struck{$x \mathrel{|o|} y \Longleftrightarrow$ ($\{x, y\} \in A$, ou $x = y$ et $x$ sommet \ill{})} \struck{des} \struck{points} \add{sommets} \struck{isolés} \struck{du graphe, \uncertain{pour} \ill{} \ill{} $(E, A)$}.
\marginal{* (\uncertain{l'ensemble} des points isolés \ill{} \uncertain{non} \ill{} $A$).}
\struck{La structure de disposition \uncertain{n'est} antiréflexive} \struck{\uncertain{équivaut} à celle de graphe simple} \add{strict}.
On trouve ainsi une bijection entre structures de disposition \uncertain{antiréflexive} sur $E$, et structures de graphe \uncertain{strict} \uncertain{sur} $E$, \uncertain{par} \uncertain{parties} de $\mathfrak{P}_2(E)$.

La structure de disposition donne une \uncertain{formulation} \uncertain{en} $\operatorname{cosupp}^{\circ}$ \ill{} \uncertain{fondamentale} des $\operatorname{cosupp}^{\circ}_{E}$ (ou $\operatorname{cosupp}^{\circ}_{\underline{E}}$, si $\underline{E} = (E, |o|_E)$) et $\operatorname{supp}^{\circ}_{E} = \operatorname{cosupp}^{\circ}_{E} \circ \operatorname{cosupp}^{\circ}_{E}$ dans $\mathfrak{P}(E)$, dont

\page{76}\note{p. 73 de l'auteur}
partie $\Sigma_{\underline{E}}$ (ou $\Sigma_E$) de $\mathfrak{P}_E$, stable par $\bigcap$ quelconques, ayant comme plus grand élément $E$, comme plus petit élément la réunion $E_0 \ldots$ \uncertain{Je} \ill{} \ill{} des Inf quelc.\ (les $\bigcap$) et Sup quelc.\ (en général $\neq \bigcup$), et \ill{} \uncertain{une} involution
\[ \complement S \overset{\mathrm{df}}{=} \operatorname{cosupp}^{\circ}_{E} A \qquad (A \in \Sigma_E \subset \mathfrak{P}(E)), \]
\uncertain{une} application canonique
\[ x \longmapsto \sigma(x) = \operatorname{supp}^{\circ}_{E}(\{x\}) : E \xrightarrow{\ \sigma\ } \Sigma_E \]
\marginal{\uncertain{On a} $x \mathrel{|o|} y \Longleftrightarrow \sigma(x) \mathrel{|o|} \sigma(y)$~; \uncertain{pour} $S, T \in \Sigma_E$ \ill{} $A \mathrel{|o|} B \Longleftrightarrow \ldots$ $A \subset \operatorname{cosupp}^{\circ}(B)$ \ill{} \ill{} \ill{} $\complement$~; \uncertain{Dans} $\mathfrak{P}(E)$ \uncertain{de} \uncertain{cette} \uncertain{structure}…}
\emph{Disposition triviale}~: formant
\[ x \mathrel{|o|} y \Longleftrightarrow x \neq y \]
\uncertain{donnée} \uncertain{correspondant} à~: $A = \emptyset$, structure de graphe discrète. On a alors
\[ \Sigma_E = \mathfrak{P}(E). \]

\textbf{NB} Pour toute partie $A$ de $E$, on a
\begin{align*}
\operatorname{cosupp}^{\circ}_{E}(A) &= \operatorname{cosupp}^{\circ}_{E}(A \cup E_0) \supset E_0, \\
\operatorname{supp}^{\circ}_{E}(A) &= \operatorname{supp}^{\circ}_{E}(A \cup E_0) \supset E_0,
\end{align*}
\note{la seconde ligne porte « $\operatorname{supp}^{\circ}_{E}(A \cup E_0)$ », le signe de réunion étant à demi formé.}
et par suite aussi
\begin{align*}
\operatorname{cosupp}^{\circ}_{E}(A) &= \operatorname{cosupp}^{\circ}(A^{*}) \\
\operatorname{supp}^{\circ}_{E}(A) &= \operatorname{supp}^{\circ}(A^{*})
\end{align*}
où
\[ A^{*} = A \setminus E_0 \cap A. \]
Donc on conclut que l'on a un iso.\ (compatible \uncertain{structure} d'ordre et $\complement$)
\[ \Sigma_E \xrightarrow{\ \sim\ } \Sigma_{E^{*}}, \qquad
\begin{array}{l} S \longmapsto S \setminus E_0 \\ S^{*} \amalg E_0 \longleftarrow S^{*} \end{array} \]
\note{la note marginale, écrite en biais sur la moitié gauche, est très serrée~; elle est rendue en partie.}

\page{77}\note{p. 74 de l'auteur}
Cela montre que pour la construction

2. \emph{Spatiales}. Dans \uncertain{ce} \uncertain{paragraphe} de la \uncertain{discussion}.

\emph{Définition 1}. Une \emph{spatiale} est un ensemble $\Sigma$, muni d'\struck{\ill{}} relations \struck{$\leq$} $\leq$, \struck{\ill{}} et d'une involution $\complement$ (relations \uncertain{aussi}, \uncertain{mais} \uncertain{aussi}~!), satisfaisant les conditions

\emph{Spat 1}. La relation $\leq$ est une \emph{relation d'ordre}, admettant des Sup et des Inf quelconques (donc un plus petit élément, noté $0_\Sigma$, et un plus grand élément, noté $1_\Sigma$).

\emph{Spat 2}. $\complement$ est une \emph{involution} ($\complement(\complement x) = x$) et « renverse la relation d'ordre »~:
\[ x \leq y \Longleftrightarrow \complement y \leq \complement x. \]
\struck{\uncertain{Spatialités}}
\[ \left[\begin{array}{ll} \text{\emph{Cor 1}} & \complement(\operatorname*{Sup}_i x_i) = \operatorname*{Inf}_i \complement x_i, \\
 & \complement(\operatorname*{Inf}_i x_i) = \operatorname*{Sup}_i \complement x_i. \\
\text{\emph{Cor.\ 2}} & \complement 0 = 1,\quad \complement 1 = 0 \end{array}\right. \]
\note{les trois lignes suivantes sont barrées de quatre traits obliques.}
\struck{À une spatiale est associée une structure de disposition, en posant}
\[ \text{\struck{$x \mathrel{|o|} y \overset{\mathrm{def}}{\Longleftrightarrow} x \leq \complement y$}} \qquad \text{\struck{($\Longleftrightarrow y \leq \complement x$ par Spat.\ 2).}} \]

\emph{Spat 3}. $\forall x \in \Sigma$, on a
\[ x \wedge \complement x = 0_\Sigma \]
ce qui équivaut par Spat 2 (Cor.\ 1 et 2) à
\[ x \vee \complement x = 1_\Sigma. \]

\page{78}\note{p. 75 de l'auteur}
À une spatiale est \struck{associée} \add{munie} d'une structure de disposition canonique~:
\[ x \mathrel{|o|} y \overset{\mathrm{def}}{\Longleftrightarrow} x \leq \complement y \qquad (\Longleftrightarrow y \leq \complement x, \text{ par Spat 2}). \]
Notons que
\[ x \mathrel{|o|} y,\quad x' \leq x,\quad y' \leq y \Longrightarrow x' \mathrel{|o|} y'. \]
\marginal{\textbf{NB} Si $x = \operatorname{Sup} x_i$, $y = \operatorname{Sup} y_j$, \uncertain{alors} $x \mathrel{|o|} y$ \ill{} $\forall i, j$, $x_i \mathrel{|o|} y_j$ \ill{}}
\struck{La structure de disposition permet de \ill{} réciproquement la structure $\leq$, $\complement$, \uncertain{puisque} \ill{}}
\note{ces deux lignes sont biffées horizontalement et barrées de quatre traits verticaux.}

\struck{\emph{Rapport}. Soit $\Sigma$ une spatiale,} considérons

\emph{Exemple de spatiales}. Si $\underline{E} = (E, |o|)$ est une disposition, $(\Sigma_{\underline{E}}, \subset, \complement)$ est une spatiale. Les \struck{relations} propriétés Spat 1, Spat 2 \uncertain{sont} \uncertain{dites} ci-dessus, elles ne dépendent que du fait que $|o|_E$ soit symétrique (Disp 1), sans plus. La propriété Spat 3 résulte de (\uncertain{Disp.\ 2}), qui implique, pour tout $A \in \mathfrak{P}(E)$
\[ A \cap \operatorname{cosupp}^{\circ}_{E}(A) \subset E_0, \]
\uncertain{en particulier} si $A \in \Sigma_E$
\[ A \cap \operatorname{cosupp}^{\circ}_{E}(A) = E_0. \]
\note{ici une longue note marginale en biais, d'une quinzaine de lignes, couvre toute la moitié inférieure gauche de la page et déborde sur le texte~; on en lit seulement~: « Ensemble ordonné \uncertain{disposé}… », « Ex.~: \ill{} », « $\Sigma_E$ », « \uncertain{Abs} », « $\operatorname{cosupp}^{\circ}$ ». Le reste est \ill{}.}
\uncertain{Donc}, si $\underline{E}$ est une disposition sur un ensemble, on \uncertain{en déduit} \struck{des dispositions} une spatiale associée \uncertain{et une} application canonique
\[ E \longrightarrow \Sigma_{\underline{E}}. \]

\page{79}\note{p. 76 de l'auteur}
La structure $\mathrel{|o|}$ sur une spatiale permet de récupérer la structure $\leq$, $\complement$ de la spatiale, grâce à la proposition suivante~:
\marginal{\uncertain{Lemme}~: \uncertain{Carac.} \uncertain{Décisive} \uncertain{des} \uncertain{relations} \uncertain{d'ordre} \ill{}~: $z \mathrel{|o|} y \Rightarrow z \mathrel{|o|} x$ \ill{} $\Leftrightarrow y \leq x$ \ill{}}

\emph{Proposition 1}. Soit $(\Sigma, \leq, \complement)$ une spatiale, $\mathrel{|o|}$ la structure de disposition associée. Alors l'application canonique
\[ \sigma : \Sigma \longrightarrow \Sigma_{\Sigma} \]
\struck{\uncertain{dispose}} est un \emph{isomorphisme de spatiales}.
L'\struck{appl} isomorphisme inverse est donné par
\[ \sigma^{-1} : S \longmapsto \operatorname{Sup} S. \]
\marginal{\textbf{NB} Le \uncertain{critère} \uncertain{des} \ill{}, \uncertain{Dans} \uncertain{le} \uncertain{cas} \ill{} $|o|$ \ill{}~: $\complement x = \operatorname{Sup}$ \ill{} $x \mathrel{|o|} \ldots$}

\emph{Scholie}. Pour un ensemble donné $\Sigma$,

(i) \uncertain{il y a} corr.\ 1--1 entre les deux ensembles suivants~:

(i) Les \uncertain{\ill{}} des structures $(\leq, \complement)$, structures de spatiales, \uncertain{sur} $\Sigma$. $\mathcal{S}$

(ii) L'ens.\ des structures de disposition $|o|$ sur $\Sigma$, telles que l'application canonique
\[ (*) \qquad \Sigma \longrightarrow \Sigma_{\underline{\Sigma}} \qquad (\text{où } \underline{\Sigma} = (\Sigma, |o|)) \]
soit bijection. $\mathcal{D}$

\struck{À l'appl.} La bijection $\mathcal{S} \xrightarrow{\ \sim\ } \mathcal{D}$ \struck{est donnée} \add{associe} \struck{par} à une structure de spat., la structure de disp.\ associée. La bijection inverse associe à une structure de disposition $\in \mathcal{D}$, la structure de spatiale sur $\Sigma$ déduite par \uncertain{transport} de str.\ des \uncertain{bijections} $(*)$.

\page{80}\note{p. 77 de l'auteur}
Ainsi, les spatiales peuvent être considérées comme des dispositions particulières, qu'on pourrait appeler les \emph{dispositions canoniques}. Si $(E, |o|)$ est une disposition quelconque, on aimerait considérer $\Sigma_E$, munie de
\[ \sigma_E : E \longrightarrow \Sigma_E, \]
comme une « disposition canonique » (ou spatiale) \emph{enveloppe} de $E$. Je voudrais dire \uncertain{ceci}~: \uncertain{à iso.\ unique près}, la spatiale $\Sigma_E$ et l'application $\sigma$.

\emph{Proposition 2}. Soient $(E, |o|)$ une disposition, $(\Sigma, \leq, \complement)$ une spatiale, et
\[ \alpha : E \longrightarrow \Sigma \]
une application. Pour qu'il existe un diagramme commutatif

\begin{tikzcd}[column sep=large]
 & \Sigma \arrow[dd, leftrightarrow, "\wr\,\tilde\alpha"] \\
E \arrow[ur, "\sigma"] \arrow[dr, "\sigma_E"'] & \\
 & \Sigma_E
\end{tikzcd}

\note{l'étiquette de la flèche oblique supérieure est écrite sur un mot raturé et se lit « $\sigma$ »~; on attendrait $\alpha$. La flèche verticale porte deux pointes, et l'étiquette « $\wr\,\tilde\alpha$ » (lecture incertaine du tilde).}
où $\tilde\alpha$ est un iso, il faut et il suffit que $\alpha$ satisfasse les deux conditions suivantes~:

a) $\forall x, y \in E$, \quad $x \mathrel{|o|} y \Longleftrightarrow$ \struck{$\alpha(x) \mathrel{|o|} \alpha(y)$}
\note{la condition a) s'interrompt au bas de la page, le membre de droite étant raturé~; la suite de la proposition~2 continue au-delà de ce lot.}

\end{document}
